\section{Análisis de Negocio}\label{sec:negocio}

Con el objetivo de validar la viabilidad estratégica de ParkLink y establecer un marco claro para la toma de decisiones, se aplicaron dos herramientas de análisis de negocios ampliamente utilizadas en la industria: El modelo de las 5C's, que permite mapear el entorno competitivo y colaborativo del proyecto, y el esquema de las 4P's, que estructura la propuesta de valor desde una perspectiva de marketing y posicionamiento.

\subsection{Análisis de las 5C's}

\subsubsection{Compañía (\textit{Company})}

ParkLink es una plataforma web y móvil que actúa como marketplace transaccional de estacionamiento urbano. Su propuesta de valor central consiste en conectar a conductores que necesitan un lugar donde estacionar con propietarios de espacios ociosos o cocheras comerciales, facilitando la reserva anticipada, la gestión de disponibilidad y el pago digital en un único entorno unificado.

La estructura del sistema se organiza en cuatro módulos de trabajo independientes (mapeo y geolocalización, reserva dinámica, autenticación y seguridad, y pagos), cuyo detalle técnico se desarrolla en la sección de Metodología. Este enfoque modular permite un desarrollo incremental alineado con los recursos disponibles del equipo.

La detección automática de disponibilidad mediante sensores IoT, predicciones mediante sistemas de IA y la integración con sistemas municipales de tránsito se identifican como evoluciones futuras fuera del alcance del MVP, dada la complejidad técnica y los recursos externos que requieren.

El principal riesgo operativo del modelo es el problema del arranque en frío: Para que la plataforma funcione, se necesitan simultáneamente propietarios que ofrezcan espacios y conductores que los demanden. La estrategia para mitigar esto consiste en lanzar en una zona geográfica acotada de CABA, donde sea posible generar masa crítica en ambos lados del mercado antes de escalar.

\subsubsection{Clientes y Segmentación (\textit{Customers})}

ParkLink atiende a cuatro segmentos de usuarios con necesidades diferenciadas:

\begin{enumerate}
    \item \textbf{Conductores particulares:} Representan el lado de la demanda. Su necesidad principal es encontrar un espacio disponible, con precio conocido de antemano, sin necesidad de circular en busca de lugar. Es el segmento más fácil de captar inicialmente porque comprende el valor del producto de forma inmediata.

    \item \textbf{Propietarios de espacios privados:} Personas con cocheras o entradas de garaje que permanecen vacías durante determinados períodos del día. Su motivación es generar ingresos pasivos a partir de un activo subutilizado, sin necesidad de gestión adicional.

    \item \textbf{Playas y cocheras comerciales:} Estacionamientos que buscan incrementar su tasa de ocupación mediante reservas anticipadas y captación de usuarios a través de un canal digital. Aportan mayor volumen de oferta y pueden acceder a herramientas de gestión de tarifas y disponibilidad.

    \item \textbf{Flotas, empresas y conductores frecuentes:} Repartidores, técnicos de campo, vehículos corporativos que necesitan previsibilidad en el estacionamiento. Representan un segmento con alta recurrencia y potencial para planes de suscripción mensual.
\end{enumerate}

Un quinto segmento de interés es el \textbf{municipio o actores públicos}, que podrían convertirse en clientes institucionales en una etapa posterior, una vez que la plataforma demuestre capacidad para gestionar datos de movilidad urbana a escala. Este segmento queda fuera del alcance del MVP por los largos procesos de licitación y los requisitos técnicos y legales que implica.

En términos de modelo de cobro, se definió que el único pagador fuera el conductor, mediante un cargo de servicio sobre cada reserva concretada, y que el propietario no tuviera ningún costo por publicar ni por cada operación. De esta forma se evita cobrar al lado de la oferta, que es el más difícil de captar en las primeras etapas de un mercado de dos lados.

\subsubsection{Competidores (\textit{Competitors})}

El análisis detallado de los competidores —incluyendo sus ventajas y limitaciones frente a ParkLink— se desarrolla en la Sección~\ref{sec:estado_arte}. A modo de síntesis, los actores relevantes son: El estacionamiento libre en la vía pública, los sistemas de estacionamiento medido (Blinkay), los navegadores con información de tráfico (Waze, Google Maps), los clasificados digitales (MercadoLibre, OLX) y las plataformas internacionales de referencia (SpotHero, EasyPark). Ninguno de los actores locales cubre de forma integral la intersección entre reserva anticipada, disponibilidad en tiempo real, pago digital y acceso a espacios privados que define la propuesta de ParkLink.

\subsubsection{Colaboradores (\textit{Collaborators})}

Para conseguir un mayor volumen de clientes, ParkLink deberá valerse de diversas alianzas estratégicas con actores clave del ecosistema de movilidad urbana y gestión de espacios en grandes ciudades:

\begin{itemize}
    \item \textbf{Propietarios y administradores de cocheras:} Su disposición a incorporarse a la plataforma define el tamaño de la oferta disponible.
    \item \textbf{Usuarios de la plataforma:} Su experiencia y satisfacción determinarán la tasa de retención y el crecimiento orgánico a través de recomendaciones.
    \item \textbf{Administradores de edificios:} Canal natural de acceso a propietarios de cocheras residenciales en zonas de alta densidad.
    \item \textbf{Municipalidades:} Fuente de información sobre regulaciones, zonas de estacionamiento medido y potenciales alianzas a futuro.
\end{itemize}

\subsubsection{Contexto del Entorno (\textit{Context})}

El entorno en el que ParkLink se desarrolla presenta condiciones favorables para su adopción:

\begin{itemize}
    \item \textbf{Social:} La población urbana ya utiliza con normalidad aplicaciones de mapas, reservas y pagos digitales. La adopción de plataformas de economía colaborativa estableció patrones de comportamiento compatibles con el modelo de ParkLink.
    \item \textbf{Tecnológico:} La penetración de smartphones en Argentina supera el 80\%, y la conectividad móvil en zonas urbanas permite el uso fluido de aplicaciones con geolocalización en tiempo real.
    \item \textbf{Económico:} En contextos de inflación, la monetización de activos ociosos como cocheras resulta especialmente atractiva para propietarios. Para conductores, conocer el precio exacto antes de llegar reduce la incertidumbre económica del traslado.
    \item \textbf{Legal y Regulatorio:} El uso de espacios privados como cocheras para alquiler temporal tiene respaldo en el Código Civil y Comercial, bajo la figura de la locación. El tratamiento de datos personales está regulado por la Ley 25.326 y la normativa de la AAIP, detalladas en el \hyperref[anexo:legal]{Anexo 3}.
\end{itemize}

\subsection{Análisis de las 4P's}

\subsubsection{Producto (\textit{Product})}

ParkLink es una aplicación móvil y plataforma web que permite encontrar, reservar y pagar espacios de estacionamiento en la Ciudad Autónoma de Buenos Aires.

El sistema contempla tres experiencias de usuario diferenciadas:

\begin{itemize}
    \item \textbf{Experiencia del conductor:} Accede a un mapa interactivo con los espacios disponibles cercanos, visualiza el precio y el horario de disponibilidad de cada uno, realiza una reserva indicando la duración estimada, y llega directamente al espacio confirmado sin necesidad de circular en busca de lugar.
    \item \textbf{Experiencia del propietario particular:} Publica su cochera o espacio libre, define los horarios en que está disponible y el precio por hora, y recibe reservas de forma automática sin necesidad de intervención adicional por cada operación.
    \item \textbf{Experiencia de playas comerciales:} Gestiona su capacidad disponible, establece tarifas diferenciadas por horario o tipo de vehículo, y recibe reservas anticipadas que le permiten optimizar la ocupación de sus instalaciones.
\end{itemize}

El diferencial competitivo central frente a los actores locales existentes es la combinación de reserva anticipada con confirmación de disponibilidad real, pago integrado y acceso a espacios privados. A diferencia de Blinkay, que informa sobre el estacionamiento medido pero no permite reservar; y a diferencia de los clasificados digitales, que facilitan alquileres mensuales pero no transacciones por horas, ParkLink resuelve el problema completo del conductor en una única interfaz.

\subsubsection{Precio (\textit{Price})}

ParkLink no cobra a los propietarios: El ingreso de la plataforma proviene exclusivamente de las transacciones exitosas y lo abona únicamente el conductor. El modelo de ingresos se estructura de la siguiente forma:

\begin{itemize}
    \item \textbf{Propietario:} Publica su cochera sin costo y recibe el 100\% del precio que fija para su espacio.
    \item \textbf{Cargo de servicio al conductor:} 8\% adicional sobre el valor de cada reserva concretada, abonado por el conductor al momento del pago como cargo de plataforma.
\end{itemize}

Este esquema garantiza que publicar sea gratuito para los propietarios y que reciban la totalidad de su tarifa, eliminando barreras de entrada al lado de la oferta. Al mismo tiempo, alinea los incentivos de la plataforma con el éxito de cada operación: ParkLink solo genera ingresos cuando hay una transacción real.

Para etapas posteriores se contempla la posibilidad de incorporar planes de suscripción mensual para cocheras comerciales con alto volumen de reservas, servicios de visibilidad destacada en el mapa, y planes para flotas empresariales con tarifas negociadas y facturación centralizada.

\subsubsection{Plaza (\textit{Place})}

La distribución de ParkLink es 100\% digital. El acceso a la plataforma se realiza a través de:

\begin{itemize}
    \item Aplicación móvil multiplataforma para iOS y Android.
    \item Plataforma web accesible desde cualquier dispositivo con navegador.
\end{itemize}

La estrategia de lanzamiento contempla un despliegue geográfico progresivo:

\begin{enumerate}
    \item \textbf{MVP (año 1):} Foco en barrios de CABA con alta densidad vehicular y demanda de estacionamiento comprobada (Palermo, Recoleta y microcentro). El objetivo es generar masa crítica local antes de escalar.
    \item \textbf{Mediano plazo (año 2):} Expansión a toda la Ciudad Autónoma de Buenos Aires y primeros municipios del Gran Buenos Aires.
    \item \textbf{Largo plazo:} Expansión a otras ciudades argentinas de alta densidad (Córdoba, Rosario, Mendoza) y evaluación de mercados de América Latina.
\end{enumerate}

\subsubsection{Promoción (\textit{Promotion})}

La estrategia de adquisición de usuarios se diferencia según el segmento objetivo:

\paragraph{Estrategia de demanda (conductores):}
\begin{itemize}
    \item Campañas en Google Ads segmentadas por búsquedas relacionadas con estacionamiento en CABA.
    \item Publicidad geolocalizada en Instagram y TikTok orientada a conductores en zonas de alta congestión.
    \item Presencia en comunidades de conductores y grupos de afinidad automotriz en redes sociales.
\end{itemize}

\paragraph{Estrategia de oferta (propietarios):}
\begin{itemize}
    \item Contacto directo con administradores de edificios de alta densidad en las zonas de lanzamiento.
    \item Campañas en redes sociales orientadas a propietarios de inmuebles en CABA.
    \item Comunicación enfocada en el ingreso pasivo: «Tu cochera vacía puede generar ingresos mientras no la usás».
\end{itemize}

\paragraph{Tácticas de crecimiento orgánico y referidos:}
\begin{itemize}
    \item Programa de referidos entre propietarios: Visibilidad destacada en el mapa por cada nuevo espacio incorporado a través de una recomendación.
    \item Visibilidad destacada en el mapa para los primeros propietarios en cada nueva zona de expansión, como incentivo de adopción temprana.
    \item Generación de contenido sobre movilidad urbana y estacionamiento en CABA, posicionando a ParkLink como referente del tema en medios digitales.
\end{itemize}
